\section{Mục tiêu của bài tập}

Theo định hướng của môn Nhập môn Trí tuệ Nhân tạo, bài tập này không tập trung vào việc xây dựng một AI tự động giải game. Thay vào đó, mục tiêu chính là xây dựng một môi trường hoạt động ổn định và áp dụng một thuật toán tìm kiếm cơ bản để xử lý cơ chế của trò chơi.

Các mục tiêu cụ thể gồm:
\begin{itemize}
    \item \textbf{Thiết kế thuật toán:} Cài đặt cơ chế flood fill bằng thuật toán DFS ở dạng lặp (sử dụng Stack thay vì đệ quy) để tránh lỗi tràn bộ nhớ và đảm bảo độ phức tạp tuyến tính O(N).
    \item \textbf{Xây dựng giao diện:} Thiết kế giao diện bằng Pygame, hỗ trợ đầy đủ thao tác chuột và ba mức độ: Dễ (5×5), Trung bình (9×9) và Khó (16×16).
    \item \textbf{Tổ chức kiến trúc:} Tách riêng phần xử lý logic của bàn chơi và phần hiển thị giao diện, giúp hệ thống dễ bảo trì, mở rộng và kiểm thử. Đồng thời, toàn bộ quá trình hoạt động của thuật toán DFS cũng được mô tả chi tiết trong báo cáo.
\end{itemize}







